\subsection{Iteration 2: Refactored Environment and Agent}
\subsubsection{About}
The second iteration involved refactoring the MiniDungeons environment to allow for streamlined integration of new levels. 
The iteration also involved the creation of a deterministic agent in order to focus on the quality of the generated levels, 
rather than the quality of the random steps that happened to be made by the included random agent.

The goal of the refactor was not to change the underlying functionality of MiniDungeons, but to reconfigure the environment for better
interoperability within the project. To achieve this, a copy of \texttt{gym-md} was created, but with all its configuration file searching removed.
Instead of using configuration files, these values would be hardcoded into the environment as it was inconsequential to the generator to have 
variable configurations of the environment.

Following this change, the renderer was updated to use pygame instead of matplotlib. Using pygame for human rendering allowed for 
improved verification of the agent behavior and level quality due to easily being able to add visual aids to the rendering window.

The final step of this iteration was to create a deterministic agent to have a consistent level tester. The agent was based on the Treasure Seeker Persona 
described in the original MiniDungeons paper \cite{holmgard2014evolvingpersonas}. The treasure seeker was chosen as it approximately resembles a completionist 
focused player. This player type was deemed a reasonable balancing point as this type of player would strive to explore the majority of a dungeon to maximize 
their reward.

\subsubsection{Evaluation}
evaluate

\subsubsection{Reflections and decisions}
reflect and decide